\documentclass[11pt]{article}
\usepackage[T1]{fontenc}
\usepackage[utf8]{inputenc}
\usepackage{lmodern}
\usepackage[a4paper,margin=27mm,headheight=15pt]{geometry}
\usepackage{amsmath,amssymb,amsthm,mathtools,bm,mathrsfs}
\usepackage{booktabs,longtable,array}
\usepackage{microtype,xcolor}
\usepackage{enumitem}
\usepackage{fancyhdr}
\usepackage[colorlinks=true,linkcolor=blue!45!black,citecolor=blue!45!black,urlcolor=blue!45!black,breaklinks=true]{hyperref}
\usepackage{bookmark}
\usepackage{listings}
\hypersetup{pdftitle={Five OEIS Asymptotic Conjectures for Distinct-Partition Norms},pdfauthor={Research report prepared for Vladimir Reshetnikov},pdfsubject={Weighted distinct partitions, asymptotic expansions, Lambert W inversion, Gumbel laws}}
\pagestyle{fancy}\fancyhf{}
\fancyhead[L]{\small Distinct-partition norms}
\fancyhead[R]{\small OEIS research report}
\fancyfoot[C]{\thepage}
\setlength{\parskip}{4pt}\setlength{\parindent}{1.2em}
\setlength{\emergencystretch}{2em}
\numberwithin{equation}{section}
\newtheorem{theorem}{Theorem}[section]
\newtheorem{proposition}[theorem]{Proposition}
\newtheorem{lemma}[theorem]{Lemma}
\newtheorem{corollary}[theorem]{Corollary}
\theoremstyle{definition}\newtheorem{definition}[theorem]{Definition}
\theoremstyle{remark}\newtheorem{remark}[theorem]{Remark}
\newcommand{\N}{\mathbb N}\newcommand{\R}{\mathbb R}\newcommand{\E}{\mathbb E}\newcommand{\Prob}{\mathbb P}
\newcommand{\dd}{\,\mathrm d}\newcommand{\ii}{\mathrm i}
\newcommand{\normp}{\operatorname{N}}
\newcommand{\coef}[2]{[#1^{#2}]}
\newcommand{\OEIS}[1]{\href{https://oeis.org/#1}{\texttt{#1}}}
\newcommand{\asym}{\mathscr H_{\alpha}}
\newcommand{\eps}{\varepsilon}
\DeclareMathOperator{\Var}{Var}
\DeclareMathOperator{\sech}{sech}
\lstset{basicstyle=\ttfamily\small,breaklines=true,columns=fullflexible,frame=single,framerule=.3pt,showstringspaces=false}
\begin{document}
\begin{titlepage}
\thispagestyle{empty}
\vspace*{12mm}
{\large\sffamily OEIS \; / \; ASYMPTOTIC ANALYSIS}\par
\vspace{11mm}
{\huge\bfseries Five OEIS Asymptotic Conjectures\par for Distinct-Partition Norms\par}
\vspace{7mm}
{\Large Proofs, all-orders logarithmic expansions,\par exact-saddle corrections, and inverse\par and extreme-value laws\par}
\vspace{13mm}
{\large Research report prepared for Vladimir Reshetnikov\par}
{\large October 1, 2026\par}
\vspace{12mm}
\hrule\vspace{6mm}
\noindent\textbf{Central family}
\[
 A_\alpha(q)=\prod_{k\ge1}(1+k^\alpha q^k),\qquad
 a_\alpha(n)=\coef{q}{n}A_\alpha(q),\qquad \alpha>0.
\]
\noindent The cases $\alpha=1,2,3,4,5$ are \OEIS{A022629},
\OEIS{A092484}, \OEIS{A265840}, \OEIS{A265841}, and \OEIS{A265842}.
All five consulted entries label their leading logarithmic asymptotic as a
conjecture. This report proves those formulas and develops a single refinement
mechanism for every fixed positive $\alpha$.
\vspace{6mm}\hrule
\vfill
\noindent\textbf{Scope and verification.} The results below have detailed English
proofs and accompanying exact and numerical checks. They are not claimed to be
Lean-verified or peer-reviewed. The literature comparison is targeted, not an
exhaustive priority certification. The complete exponentially improved
coefficient transseries is not claimed: the report distinguishes its proved
logarithmic block, exact-saddle expansion, and Fourier-resonance scale.
\end{titlepage}

\begin{abstract}
For a partition $\lambda$ into distinct parts, let $\normp(\lambda)$ be the
product of its parts. We analyze the fixed positive norm moments
$a_\alpha(n)=\sum_{\lambda\vdash n\text{ distinct}}\normp(\lambda)^\alpha$.
A direct argument proves the logarithmic growth formulas still marked
conjectural in five OEIS entries. With $K=\sqrt{2n}$, $r=\log K$, and
$c=\pi^2/(6\alpha^2)$, we then prove, to every fixed order,
\[
 \log a_\alpha(n)=\alpha K\left(r-1+\frac c r+\frac c{r^2}
 +\frac{c-c^2/2}{r^3}+\cdots\right).
\]
An explicit inverse-function construction generates all coefficients; seven are
printed. A separate local-limit argument gives multiplicatively accurate
exact-saddle asymptotics and all fixed Edgeworth orders. This separation matters:
exponentiating a fixed truncation of the logarithmic expansion does not give a
relative asymptotic equivalent. We invert the dominant block using the principal
Lambert function and obtain asymptotics for the first index attaining a large
value. Under the norm-biased measure, the occupation boundary has a logistic
profile, while the largest part has a Gumbel limit and is asymptotic to
$(1+1/\alpha)\sqrt{2n}$. Finally, we prove that the first Fourier resonances have
logarithmic modulus of order $-K/(\alpha^3(\log K)^3)$, explaining a striking
finite-size limitation visible in the exact computations. Reproducible scripts,
source-status notes, and further research questions accompany the proofs.
\end{abstract}
\setcounter{tocdepth}{1}
\begingroup
\setlength{\parskip}{0pt}
\tableofcontents
\endgroup
\newpage

\section{The problem, its significance, and the source audit}
\subsection{Five conjectures in one natural family}
Let $\mathcal D_n$ be the set of partitions of $n$ into distinct positive parts.
The empty partition has norm $1$. For fixed $\alpha>0$, define
\begin{equation}\label{eq:def}
 \normp(\lambda)=\prod_{k\in\lambda}k,
 \qquad a_\alpha(n)=\sum_{\lambda\in\mathcal D_n}\normp(\lambda)^\alpha.
\end{equation}
Choosing whether each part occurs gives the locally uniformly convergent product
\begin{equation}\label{eq:product}
 A_\alpha(q)=\sum_{n\ge0}a_\alpha(n)q^n
 =\prod_{k\ge1}(1+k^\alpha q^k),\qquad |q|<1.
\end{equation}
For integral $\alpha$ the coefficients are integers. One may interpret $k^\alpha$
as the number of colors available for a part of size $k$, with at most one part
of each \emph{size}, regardless of color. This is not the product
$\prod_k(1+q^k)^{k^\alpha}$, which allows several differently colored copies of
the same size.

The OEIS entries consulted on October 1, 2026 have the following status
\cite{oeis1,oeis2,oeis3,oeis4,oeis5}.
\begin{center}
\begin{tabular}{cll}
\toprule
$\alpha$ & Sequence & Date of the conjectural formula\\
\midrule
1 & \OEIS{A022629} & May 8, 2018\\
2 & \OEIS{A092484} & December 27, 2020\\
3 & \OEIS{A265840} & December 27, 2020\\
4 & \OEIS{A265841} & December 27, 2020\\
5 & \OEIS{A265842} & December 27, 2020\\
\bottomrule
\end{tabular}
\end{center}
All five formulas are attributed there to V\'aclav Kote\v{s}ovec and have the
unified form
\begin{equation}\label{eq:conjecture}
 \log a_\alpha(n)\sim
 \alpha\sqrt{\frac n2}\bigl(\log(2n)-2\bigr).
\end{equation}
The full array of nonnegative integral moments is \OEIS{A292189}; its columns
$0,1,\ldots,5$ are identified in that entry \cite{oeisarray}. The zeroth column
is the ordinary distinct-partition function, but $\alpha=0$ is a different
asymptotic regime and is excluded from our fixed-positive-parameter theorems.

These sequences are not arbitrary numerical constructions. They simultaneously
encode weighted enumeration, moments of a fundamental multiplicative partition
statistic, and a conditioned model of independent binary occupancies. Their
common structure makes a family theorem more informative than five separate
fits to numerical data.

\subsection{What is established here, and what is not a priority claim}
The conjectural labels in the consulted entries are verifiable facts; the claim
that no proof has ever appeared anywhere would be much stronger and is not made.
This report supplies self-contained proofs of the displayed formulas, followed
by explicit refinements and probabilistic consequences. Novelty of the
refinements should be assessed against the broader literature before publication.

Schneider and Sills survey partition norms, including extremal norms over
distinct partitions and their connection with triangular numbers
\cite{schneider}. Bridges and Craig study norm distributions and moments for
\emph{unrestricted} partitions \cite{bridges}; repeated parts change the
analytic problem substantially. Rana, Kaur, and Kumar give recurrences and
related norm identities for distinct partitions \cite{rana}. None of those
results is used as an unproved asymptotic input here.

A particularly relevant recent source is Gikunda's March 2026 dissertation
\cite{gikunda}. Its Chapter 4 uses
$\prod_k(1+u^{\log k}q^k)^{b_k}$ and proves log-norm central limit theorems.
For $b_k=1$ and $u=e^\alpha$ this is our generating family. However, the
coefficient estimate in Proposition 4.2.5 and the surrounding analysis use a
\emph{shrinking} tilt $\log u=O(t^\delta)$ as the saddle parameter $t$ tends
to zero. Our tilt $\alpha>0$ stays fixed. Thus that estimate does not directly
supply the fixed-tilt results below. This is a distinction between asymptotic
regimes, not a claim that the underlying family has not previously been studied.

\subsection{Relation to the ProveIt repository}
The repository snapshot consulted was
\texttt{1f1981f682b2878bde51a6ad40c22777f362fc05}.
Its \texttt{Analysis/Transseries/README.md} describes the transseries and
inversion volumes, accompanying verification programs, and the separate Lean
library for asymptotic scales and power-logarithmic monomials \cite{proveit}.
The present report follows that program by separating three precision levels:
a logarithmic asymptotic block, a relative exact-saddle expansion, and flat
Fourier contributions. Its Lambert-$W$ inversion is given explicitly rather
than left as a formal instruction to ``reverse the series.''

No repository theorem is assumed here without proof. In particular, the
existence of general transseries infrastructure is not treated as a formal
certificate for our analytic remainder estimates. The new material is an
independent, reproducible research contribution suitable for integration after
review; it does not modify the repository.

\section{Main results and the hierarchy of scales}
Throughout the paper, $\alpha>0$ is fixed. Constants in $O$-terms may depend on
$\alpha$ and on the indicated truncation order. Put
\begin{equation}\label{eq:Kr}
 K=\sqrt{2n},\qquad r=\log K,\qquad c=\frac{\pi^2}{6\alpha^2}.
\end{equation}

\begin{theorem}[All-orders logarithmic expansion]\label{thm:log}
There are explicit polynomials $P_j(c)\in\mathbb Q[c]$ such that, for every
fixed integer $R\ge1$,
\begin{equation}\label{eq:logall}
 \log a_\alpha(n)=\alpha K\left(r-1+
       \sum_{j=1}^R\frac{P_j(c)}{r^j}
       +O_{\alpha,R}(r^{-R-1})\right).
\end{equation}
The first seven are
\begin{align}
 P_1(c)&=c,&P_2(c)&=c,\notag\\
 P_3(c)&=c-\tfrac12c^2,&
 P_4(c)&=c+\tfrac9{10}c^2,\notag\\
 P_5(c)&=c+\tfrac{41}5c^2+\tfrac12c^3,\label{eq:P}\\
 P_6(c)&=c+27c^2+\tfrac{1973}{210}c^3,\notag\\
 P_7(c)&=c+\tfrac{129}2c^2+\tfrac{8109}{70}c^3-\tfrac58c^4.\notag
\end{align}
Equations \eqref{eq:dformal}--\eqref{eq:formalH} below generate every $P_j$.
\end{theorem}
In particular,
\begin{align}\label{eq:firstterms}
 \log a_\alpha(n)
 ={}&\alpha K(r-1)+\frac{\pi^2K}{6\alpha r}
     +\frac{\pi^2K}{6\alpha r^2}\notag\\
 &+\left(\frac{\pi^2}{6\alpha}-\frac{\pi^4}{72\alpha^3}\right)
       \frac K{r^3}+O_\alpha(K/r^4).
\end{align}
This proves \eqref{eq:conjecture} for all five OEIS sequences, and more strongly
identifies the $-\alpha K$ term with an error $O_\alpha(K/r)$.

For a different level of precision, define
\begin{equation}\label{eq:F}
 F(t)=\sum_{k\ge1}\log(1+k^\alpha e^{-tk}),\qquad t>0,
\end{equation}
and let $t=t_n$ be the unique solution of $-F'(t)=n$. Set
\begin{equation}\label{eq:phase}
 s=-W_{-1}(-t/\alpha),\qquad M=e^s,
 \qquad t=\alpha s/M,\qquad H=\frac M{\alpha(s-1)}.
\end{equation}
These definitions use the large real branch $s>1$, valid for all sufficiently
large $n$. Let $B=F''(t)$ and $\kappa_j=(-1)^jF^{(j)}(t)$.

\begin{theorem}[Exact-saddle coefficient expansion]\label{thm:saddle}
As $n\to\infty$,
\[
 M\sim K,\quad s\sim r,\quad B\sim\frac{M^3}{\alpha s},
 \quad
 a_\alpha(n)\sim\frac{\exp(F(t_n)+nt_n)}{\sqrt{2\pi B}}.
\]
More precisely, for every fixed $J\ge1$,
\begin{equation}\label{eq:edgeworth}
 a_\alpha(n)=\frac{\exp(F(t_n)+nt_n)}{\sqrt{2\pi B}}
 \left(1+\sum_{\ell=1}^{J-1}\mathcal E_\ell
       +O_{\alpha,J}((s/M)^J)\right),
\end{equation}
where each $\mathcal E_\ell=O_{\alpha,\ell}((s/M)^\ell)$ is given explicitly
in \eqref{eq:edgepoly}. The first correction is
\begin{equation}\label{eq:E1}
 \mathcal E_1=\frac{\kappa_4}{8B^2}
                    -\frac{5\kappa_3^2}{24B^3}.
\end{equation}
\end{theorem}

\begin{theorem}[Asymptotic inversion]\label{thm:inverse}
For $Y\to\infty$, let $N_\alpha(Y)$ be the least integer $n$ such that
$a_\alpha(n)\ge e^Y$. Define
\begin{equation}\label{eq:Wlead}
 w=W_0\left(\frac{Y}{\alpha e}\right),\qquad R_0=1+w,
 \qquad K_0=\frac{Y}{\alpha w}.
\end{equation}
Then
\begin{equation}\label{eq:inverselead}
 N_\alpha(Y)=\frac{K_0^2}{2}\left(
 1-\frac{2c}{R_0^2}-\frac{2c}{R_0^3}
 +\frac{4c^2-2c}{R_0^4}+O_\alpha(R_0^{-5})\right).
\end{equation}
A formal recurrence in Section~\ref{sec:inverse} gives all inverse-logarithmic
orders, with a remainder at every fixed order.
\end{theorem}

\begin{theorem}[Logistic boundary and Gumbel extreme]\label{thm:geometry}
Give $\mathcal D_n$ the probability law
$\Prob_n(\lambda)=\normp(\lambda)^\alpha/a_\alpha(n)$.
For each fixed $x\in\R$, a part near $M+xH$ has limiting occupation probability
$(1+e^x)^{-1}$; any fixed finite collection of distinct such sites is
asymptotically independent. If $L_n$ is the largest part, define
\begin{equation}\label{eq:bdef}
 b=b(t_n)>\alpha/t_n,\qquad
 \frac{b^\alpha e^{-t_nb}}{t_n}=1.
\end{equation}
Then, for every fixed real $x$,
\begin{equation}\label{eq:gumbel}
 \Prob_n\{t_n(L_n-b)\le x\}\longrightarrow \exp(-e^{-x}),
 \qquad \frac{L_n}{\sqrt{2n}}\longrightarrow 1+\frac1\alpha
 \quad\text{in probability}.
\end{equation}
\end{theorem}

\paragraph{A precision warning.}
The error in \eqref{eq:logall} is $O(K/r^{R+1})$, which still grows without
bound for fixed $R$. Consequently, exponentiating a fixed truncation need not
approximate $a_\alpha(n)$ to relative error $o(1)$. Formula
\eqref{eq:edgeworth} does, because it retains the \emph{exact} saddle and the
\emph{exact} infinite sums. There is also a smaller exponential scale from
noncentral Fourier arcs; Section~\ref{sec:resonance} identifies it. We do not
conflate these three statements into an unjustified single explicit
coefficient transseries.

\section{A short proof of all five posted conjectures}
This section already proves the formulas that are explicitly conjectural in the
OEIS entries. Later sections supply the stronger refinements.

\begin{lemma}[A useful continuum estimate]\label{lem:rough}
For $s\to\infty$, $M=e^s$, and $t=\alpha s/M$,
\begin{equation}\label{eq:roughF}
 F(t)=\alpha M(s/2-1)+O_\alpha(M/s).
\end{equation}
\end{lemma}
\begin{proof}
The function $f_t(x)=\log(1+x^\alpha e^{-tx})$ is nonnegative and unimodal
on $(0,\infty)$, tends to zero at both ends, and has maximum $O_\alpha(s)$.
For a nonnegative unimodal function, comparison on unit intervals gives
\[
 \left|\sum_{k\ge1}f_t(k)-\int_0^\infty f_t(x)\dd x\right|
 \le 2\sup_x f_t(x)=O_\alpha(s).
\]
With $x=Mu$, the exponent is $\alpha[s(1-u)+\log u]$. Replacing
$\log(1+e^z)$ by $\max(z,0)$ changes the integral by $O_\alpha(M/s)$:
in a fixed neighborhood of $u=1$, the derivative of $s(u-1)-\log u$
is comparable with $s$, and
$\int_\R\log(1+e^{-|v|})\dd v<\infty$.
Outside that neighborhood, split off $0<u<e^{-s/2}$; its integral is
exponentially small times a polynomial in $s$. On the remaining region
away from $1$, the absolute exponent grows linearly in $s$, except for the
negligible small lower crossing. Thus the remaining error is smaller than
$M/s$. The lower crossing changes the following signed baseline by $O(1)$:
\[
 \alpha M\int_0^1\bigl(s(1-u)+\log u\bigr)\dd u
       =\alpha M(s/2-1).
\]
The lattice error $O(s)$ is also $O(M/s)$.
\end{proof}

\begin{proposition}\label{prop:elementary}
For every fixed $\alpha>0$,
\begin{equation}\label{eq:twoterms}
 \log a_\alpha(n)=\alpha\sqrt{2n}\bigl(\log\sqrt{2n}-1\bigr)
                   +O_\alpha\left(\frac{\sqrt n}{\log n}\right).
\end{equation}
\end{proposition}
\begin{proof}
Let $m$ be maximal with $m(m+1)/2\le n$, and write
$d=n-m(m+1)/2$, so $0\le d\le m$. The distinct partition
\[
 1+2+\cdots+(m-1)+(m+d)
\]
has norm at least $m!$. Therefore
\[
 \log a_\alpha(n)\ge\alpha\log(m!)
   =\alpha K(r-1)+O_\alpha(r),
\]
where $m=K+O(1)$ and the usual integral bounds, or Stirling's formula, give
$\log(m!)=m\log m-m+O(\log m)$.

For the upper bound, nonnegative coefficients give
$a_\alpha(n)\le\exp(F(t)+nt)$ for every $t>0$.
Choose $t=\alpha r/K$, that is, $M=K$ and $s=r$ in
Lemma~\ref{lem:rough}. Since $n=K^2/2$,
\[
 F(t)+nt=\alpha K(r/2-1)+\alpha Kr/2+O_\alpha(K/r)
        =\alpha K(r-1)+O_\alpha(K/r).
\]
The two bounds prove \eqref{eq:twoterms}.
\end{proof}

\begin{remark}
The lower-bound partition need not maximize the norm. An exact maximizer is
unnecessary. Also, a bare asymptotic equivalence to $\alpha K(r-1)$ would not
by itself certify $-\alpha K$ as the second term; the error in
\eqref{eq:twoterms} does certify it.
\end{remark}

\section{The exact probabilistic saddle and its scales}\label{sec:saddle}
For fixed $t>0$, let $X_k$ be independent Bernoulli random variables with
\begin{equation}\label{eq:pk}
 p_k=\Prob_t(X_k=1)=\frac{k^\alpha e^{-tk}}{1+k^\alpha e^{-tk}},
 \qquad S=\sum_{k\ge1}kX_k.
\end{equation}
The sum is finite almost surely because $\sum p_k<\infty$. Multiplying the
Bernoulli probabilities gives the exact identity
\begin{equation}\label{eq:tiltidentity}
 \Prob_t(S=n)=a_\alpha(n)e^{-nt-F(t)}.
\end{equation}
Differentiation under the absolutely convergent sums shows
\[
 \E_tS=-F'(t)=\sum k p_k,\qquad
 \Var_t S=F''(t)=\sum k^2p_k(1-p_k)>0.
\]
The mean is continuous and strictly decreasing from $\infty$ to $0$ as $t$
increases from $0$ to $\infty$. Hence the saddle equation has a unique solution
for every $n>0$.

The large crossing $k^\alpha e^{-tk}=1$ occurs at $k=M$ in
\eqref{eq:phase}. Its logarithmic derivative there is
\[
 \frac\alpha M-t=-\frac{\alpha(s-1)}M=-H^{-1}.
\]
Thus $H$ is the width of the transition from predominantly occupied to
predominantly unoccupied sites.

\begin{lemma}[Transition estimates]\label{lem:scales}
For fixed $\alpha>0$, as $s\to\infty$,
\begin{align}
 \sum_{k\ge1}kp_k&\sim M^2/2,\label{eq:meanleading}\\
 \sum_{k\ge1}k^2p_k(1-p_k)&\sim M^2H
                    \sim M^3/(\alpha s).\label{eq:varianceleading}
\end{align}
For each fixed integer $j\ge2$,
\begin{equation}\label{eq:mombound}
 \sum_{k\ge1}k^j p_k(1-p_k)=O_{\alpha,j}(M^{j+1}/s).
\end{equation}
Moreover, for every fixed nonnegative integer $h$,
\begin{equation}\label{eq:profilemoment}
 \frac1{H^{h+1}}\sum_{k\ge1}|k-M|^h p_k(1-p_k)
 \longrightarrow
 \int_\R |x|^h\frac{e^x}{(1+e^x)^2}\dd x.
\end{equation}
\end{lemma}
\begin{proof}
For $k=M+xH$ with $x$ in a fixed bounded interval,
\[
 \alpha\log k-tk
 =\alpha\log(1+xH/M)-\alpha s xH/M
 =-x+O_{\alpha,x}(s^{-2}).
\]
Here and below rounding $k$ to an integer changes the scaled coordinate by
$O(H^{-1})=o(1)$. Thus $p_k$ converges to $(1+e^x)^{-1}$.
Since $H\to\infty$, sums over any fixed $x$-interval are Riemann sums.
The limiting variance density has integral $1$.

We record a tail justification so that this argument is not just a local
heuristic. On $k/M\in[1/2,3/2]$, the derivative of the log odds is
negative and comparable to $-1/H$ for large $s$; the variance density is
therefore bounded by $C_\alpha e^{-c_\alpha|k-M|/H}$.
On $k\ge3M/2$, the odds themselves are exponentially small in $s$ and then
in $tk$. On $k\le M/2$, use
$p_k(1-p_k)\le\min(1,k^{-\alpha}e^{tk})$, and split at $M/s$.
The range $M/s<k\le M/2$ is exponentially small in $s$ after scaling by
any of the powers of $M$ above. On $k\le M/s$, the crude bound is
$O_\alpha(k^{-\alpha})$; its contribution to \eqref{eq:profilemoment},
even with $|k-M|^h\le M^h$, is $o(H^{h+1})$.
Indeed the resulting ratio is bounded by a fixed power of $s$ times
$M^{-\min(\alpha,1)}$, with an extra logarithm when $\alpha=1$.
These bounds establish uniform integrability for every fixed moment.
They also prove \eqref{eq:mombound} and \eqref{eq:varianceleading}.

Finally, $p_{\lfloor Mu\rfloor}\to1$ for $0<u<1$ and tends to zero for
$u>1$. The same tail bounds and a Riemann sum give
$M^{-2}\sum k p_k\to\int_0^1u\dd u=1/2$.
\end{proof}

At the exact saddle, \eqref{eq:meanleading} gives $M\sim K$ and $s\sim r$.
For cumulants, write $b_1(p)=p$ and
$b_{j+1}(p)=p(1-p)b_j'(p)$. Then
\begin{equation}\label{eq:cumulants}
 \kappa_j=\sum_{k\ge1}k^j b_j(p_k),\qquad
 |b_j(p)|\le C_jp(1-p)\quad(j\ge2).
\end{equation}
Consequently, with $\eps=s/M$,
\begin{equation}\label{eq:normcumulants}
 \frac{|\kappa_j|}{B^{j/2}}
      =O_{\alpha,j}(\eps^{(j-2)/2}),\qquad j\ge3.
\end{equation}
These bounds are deliberately conservative; signed cancellations can make
individual corrections smaller. They suffice for all fixed Edgeworth orders.

\section{The local limit theorem, including noncentral arcs}\label{sec:llt}
\subsection{An elementary consecutive-block inequality}
\begin{lemma}\label{lem:block}
There is an absolute constant $c_0>0$ such that for every integer $h\ge2$,
every integer $a$, and every $|\theta|\le\pi$,
\begin{equation}\label{eq:block}
 \sum_{k=a}^{a+h-1}\sin^2(k\theta/2)
       \ge c_0\min(h,h^3\theta^2).
\end{equation}
\end{lemma}
\begin{proof}
The geometric-sum formula gives
\[
 \sum_{k=a}^{a+h-1}\sin^2(k\theta/2)
 \ge\frac12\left(h-
       \left|\frac{\sin(h\theta/2)}{\sin(\theta/2)}\right|\right).
\]
When $h|\theta|\le1$, Taylor bounds for sine show that the difference in
parentheses is at least $c h(h^2-1)\theta^2$.
For $h|\theta|\ge2\pi$, the bound
$|\sin(\theta/2)|\ge|\theta|/\pi$ makes the ratio at most $h/2$.
On the intermediate range $1\le h|\theta|\le2\pi$, its ratio to $h$ is
uniformly less than one. To see the uniformity, let $h\to\infty$ and use
uniform convergence to $|\sin(x/2)/(x/2)|<1$ on $1\le x\le2\pi$;
the finitely many remaining $h$ have the same strict inequality by the
strict triangle inequality for a nonconstant geometric progression.
Taking the minimum of the resulting positive constants proves the lemma.
\end{proof}

\subsection{Fourier decay and the central Gaussian}
Let $\phi_t(\theta)=\E_t e^{\ii\theta(S-\E_tS)}$. The modulus identity
for one Bernoulli factor implies
\begin{equation}\label{eq:phibound}
 |\phi_t(\theta)|\le
 \exp\left(-2\sum_{k\ge1}p_k(1-p_k)\sin^2(k\theta/2)\right).
\end{equation}
In a consecutive block of length $h\asymp_\alpha M/s$ centered at $M$,
all $p_k(1-p_k)$ are bounded below by a positive constant. On
$|\theta|\le c_1/M$, this block gives
\begin{equation}\label{eq:centralbound}
 |\phi_t(\theta)|\le e^{-c_2 B\theta^2}.
\end{equation}
Indeed $k\asymp M$ on the block, and $\sin(k\theta/2)$ is comparable with
$k\theta$ for sufficiently small $c_1$.
On its complement in $[-\pi,\pi]$, Lemma~\ref{lem:block} gives
\begin{equation}\label{eq:noncentral}
 |\phi_t(\theta)|\le
 \exp(-c_3M/s^3),\qquad c_1/M\le|\theta|\le\pi.
\end{equation}
All constants here may depend on $\alpha$. Notice the power $s^3$, rather
than $s$: the random part sizes lie in a relatively narrow consecutive block,
so one must control possible near-resonances. An argument checking only a
neighborhood of zero would leave a genuine gap.

At the saddle, Fourier inversion gives
\[
 \Prob_t(S=n)=\frac1{2\pi}\int_{-\pi}^{\pi}\phi_t(\theta)\dd\theta.
\]
For bounded $x=\theta\sqrt B$, Taylor expansion and
\eqref{eq:normcumulants} give
$\phi_t(x/\sqrt B)\to e^{-x^2/2}$.
The central bound dominates this integrand by $e^{-c x^2}$.
The noncentral integral, even after multiplication by $\sqrt B$, is
$O(\sqrt B e^{-c_3M/s^3})=o(1)$. Thus
\begin{equation}\label{eq:llt}
 \Prob_t(S=n)\sim(2\pi B)^{-1/2}.
\end{equation}
Together with \eqref{eq:tiltidentity}, this proves the leading part of
Theorem~\ref{thm:saddle}. In particular,
\begin{equation}\label{eq:logsaddle}
 \log a_\alpha(n)=F(t_n)+nt_n+O_\alpha(s),
\end{equation}
since $\log B=O(s)$.

\subsection{A uniform variant used for conditioning}\label{sec:uniform}
The same proof gives the following useful extension. Delete either a fixed
number of sites $k=O(M)$, or all sites above a cutoff $m$ for which
$m/M\to\gamma>1$ and the deleted mean and variance are $O(M)$ and $O(M^2)$.
The central consecutive block remains available, the retained variance is
$B(1+o(1))$, and the cumulant bounds remain valid. If an integer target differs
from the retained mean by $o(\sqrt B)$, its point probability is
\begin{equation}\label{eq:uniformllt}
 (2\pi B)^{-1/2}(1+o(1)).
\end{equation}
This follows in the central integral by inserting
$e^{-\ii x(m-\E S)/\sqrt B}\to1$; all the dominating estimates are unchanged.
In the fixed-site deletion case, the deleted variance is automatically
$O(M^2)$. The assertion is uniform for boundedly many deletions and shifts
bounded by a fixed multiple of $M$.

\section{The all-orders coefficient correction}\label{sec:edgeworth}
The normalized Fourier integrand has the expansion
\begin{equation}\label{eq:centralseries}
 \phi_t(x/\sqrt B)
 =e^{-x^2/2}\exp\left(
       \sum_{j\ge3}\frac{\kappa_j(\ii x)^j}{j!B^{j/2}}\right).
\end{equation}
Assign weight $j-2$ to the $j$th normalized cumulant. Every monomial of odd
weight has odd total power of $x$ and integrates to zero. A monomial of weight
$2\ell$ has order $O(\eps^\ell)$ by \eqref{eq:normcumulants}.
Gaussian integration therefore gives
\begin{equation}\label{eq:edgepoly}
 \mathcal E_\ell=
 \sum_{\substack{m_3,\ldots,m_{2\ell+2}\ge0\\
                  \sum_{j=3}^{2\ell+2}(j-2)m_j=2\ell}}
 \frac{(-1)^{T/2}(T-1)!!}{B^{T/2}}
 \prod_{j=3}^{2\ell+2}
       \frac{\kappa_j^{m_j}}{(j!)^{m_j}m_j!},
 \qquad T=\sum_{j=3}^{2\ell+2}jm_j.
\end{equation}
Here $T$ is automatically even. Besides \eqref{eq:E1}, the next term is
\begin{align}\label{eq:E2}
 \mathcal E_2={}&-\frac{\kappa_6}{48B^3}
 +\frac{7\kappa_3\kappa_5}{48B^4}
 +\frac{35\kappa_4^2}{384B^4}\notag\\
 &-\frac{35\kappa_3^2\kappa_4}{64B^5}
 +\frac{385\kappa_3^4}{1152B^6}.
\end{align}

For completeness, the formal expansion above has the needed integrated
remainder. On $|\theta|\le c/M$, derivatives of order $j\ge2$ of
$\log\E_t e^{\ii\theta S}$ have absolute value at most
$C_{\alpha,j}M^{j+1}/s$. For $k\le2M$, take $c$ sufficiently small so that
the Bernoulli denominators have modulus comparable with $1+k^\alpha e^{-tk}$.
For $k>2M$, the odds tend uniformly to zero; the denominators are bounded
away from zero and the exponentially decreasing tails give the same bound.
These observations justify Taylor's formula with controlled remainders.

Fix $J$, choose a small $\delta>0$ depending on $J$, and first restrict the
normalized integral to $|x|\le\eps^{-\delta}$. Expand through weight $2J-1$,
using Taylor's formula for the exponential. The terms of weight $2J$ and
higher have an integrable bound of the form
$C_J\eps^J(1+|x|^{C_J})e^{-x^2/4}$, after choosing $\delta$ small enough.
The omitted odd-weight terms integrate to zero on the symmetric interval.
Between this interval and $c\sqrt B/M$, \eqref{eq:centralbound} makes the
integral smaller than every power of $\eps$. The remaining Fourier range is
bounded by \eqref{eq:noncentral}, with the same conclusion. This proves
\eqref{eq:edgeworth} and its error for every fixed $J$.

\begin{remark}[What an effective finite-size bound would contain]\label{rem:effective}
Before absorbing flat terms into the asymptotic $O$-term, the proof contains
an additional relative bound of the shape
$C_\alpha\sqrt B\exp(-c_\alpha M/s^3)$ for noncentral arcs.
This can be much larger than the central Edgeworth corrections at attainable
$n$, especially for larger $\alpha$. The theorem is not a claim that its
first few terms have small relative error uniformly in $\alpha$ or for all
moderate $n$.
\end{remark}

\section{Continuum smoothing to every logarithmic order}\label{sec:sommerfeld}
\subsection{The inverse at the occupation boundary}
Define
\begin{equation}\label{eq:hs}
 h_s(u)=s(u-1)-\log u.
\end{equation}
Near $u=1$ it is strictly increasing. Let $u_s(v)$ be its inverse with
$u_s(0)=1$, and put
\begin{equation}\label{eq:U}
 U_m(s)=u_s^{(m)}(0).
\end{equation}
The derivative operator along the inverse is
\begin{equation}\label{eq:Urec}
 \frac\dd{\dd v}=\frac1{s-1/u}\frac\dd{\dd u}.
\end{equation}
Starting with $u$ and applying \eqref{eq:Urec} repeatedly, then setting $u=1$,
gives an effective rational-function recursion. Equivalently, Lagrange
inversion gives
\begin{equation}\label{eq:Ulagrange}
 U_m(s)=(m-1)!\,[z^{m-1}]
       \left(\frac{z}{sz-\log(1+z)}\right)^m.
\end{equation}
The first required odd derivatives are
\begin{align}\label{eq:Uexplicit}
 U_1(s)&=\frac1{s-1},\notag\\
 U_3(s)&=\frac{2s+1}{(s-1)^5},\notag\\
 U_5(s)&=\frac{24s^3+58s^2+22s+1}{(s-1)^9},\\
 U_7(s)&=\frac{720s^5+3708s^4+4400s^3+1452s^2+114s+1}
                    {(s-1)^{13}}.\notag
\end{align}

Write $\eta(z)=(1-2^{1-z})\zeta(z)$. The elementary even kernel
$g(v)=\log(1+e^{-|v|})$ has moments
\begin{equation}\label{eq:kernelmom}
 \int_\R v^{2j}g(v)\dd v=2(2j)!\eta(2j+2),\qquad j\ge0.
\end{equation}
Indeed expand the logarithm on $(0,\infty)$ and integrate each exponentially
decaying term. Odd moments vanish.

\begin{lemma}[Boundary smoothing expansion]\label{lem:smoothing}
For every fixed $J\ge1$,
\begin{align}\label{eq:smoothing}
 \int_0^\infty\log(1+x^\alpha e^{-tx})\dd x
 =M\left(\alpha(s/2-1)+
      \sum_{j=1}^J\frac{2\eta(2j)}{\alpha^{2j-1}}U_{2j-1}(s)
      +O_{\alpha,J}(s^{-2J-1})\right),
\end{align}
where $t=\alpha s e^{-s}$. The expansion may be differentiated with respect
to $s$ any fixed number of times, retaining a remainder smaller than any
prescribed inverse power of $s$ by taking sufficiently many terms first.
\end{lemma}
\begin{proof}
After $x=Mu$, the integrand is $\log(1+e^{-\alpha h_s(u)})$.
The baseline is
$\alpha M\int_0^1[s(1-u)+\log u]\dd u=\alpha M(s/2-1)$.
The discrepancy between this signed baseline and the positive-part baseline
is confined to the small lower crossing and is exponentially small in $s$
before multiplication by $M$.

Choose a fixed small $\rho>0$ and restrict initially to
$|u-1|<\rho$. Replacing the original integrand by its positive-part baseline
leaves $g(\alpha h_s(u))$. On a symmetric $v$-interval $|v|<\gamma s$
inside the image of this neighborhood, changing variables $v=h_s(u)$ gives
\[
 \int g(\alpha v)u_s'(v)\dd v.
\]
The omitted portions have exponentially small tails. Taylor-expand
$u_s'(v)$ through degree $2J-1$. On the chosen interval, repeated use of
\eqref{eq:Urec} bounds the derivative of order $2J$ of $u_s'$ by
$O_J(s^{-2J-1})$. The Taylor remainder is thus bounded by
$C_Js^{-2J-1}|v|^{2J}$, whose product with $g(\alpha v)$ is integrable.
After extending the polynomial terms to the whole real line, odd terms
vanish, and \eqref{eq:kernelmom} supplies exactly
$2\eta(2j)U_{2j-1}(s)/\alpha^{2j-1}$.

Here is a more explicit treatment of the region near zero, which must not
be mistaken for part of the main boundary. On $0<u<e^{-s/2}$, both the
original integrand and the absolute baseline have integral at most
$C_\alpha s e^{-s/2}$. On
$e^{-s/2}\le u\le1-\rho$, the concavity of
$s(1-u)+\log u$ puts it above a positive multiple of $s$ for large $s$.
The smoothing error there is exponentially small. On $u\ge1+\rho$ the
log odds are negative and decrease at rate comparable with $s$.
All these contributions are $O_{\alpha,J}(s^{-2J-1})$ for every fixed $J$.

Differentiation with respect to $s$ inserts powers of $1-u$ and derivatives
of the same logistic kernel. On the boundary interval the inverse derivatives
still have polynomial bounds in $s^{-1}$; away from it, exponential tail bounds
remain after multiplication by any fixed polynomial. Taylor expansion to a
higher order therefore gives the differentiated version with any prescribed
inverse-power precision. No assumption about differentiating an arbitrary
$O$-term is being used.
\end{proof}

\subsection{From the continuum to the lattice}
As in Lemma~\ref{lem:rough}, the difference between $F(t)$ and its continuum
integral is $O_\alpha(s)$. We also need the mean. The function
\[
 x\longmapsto \frac{x^{\alpha+1}e^{-tx}}{1+x^\alpha e^{-tx}}=xp(x)
\]
is unimodal and has maximum $O_\alpha(M)$. To verify unimodality, its derivative
has the sign of $1+(\alpha-tx)(1-p(x))$. This is positive for
$x\le\alpha/t$. For $x>\alpha/t$, the function
$(tx-\alpha)(1-p(x))$ increases strictly from zero to infinity. Thus
\begin{equation}\label{eq:meanlattice}
 \sum_{k\ge1}kp_k=\int_0^\infty xp(x)\dd x+O_\alpha(M).
\end{equation}
This exact lattice error is negligible at every fixed inverse-logarithmic
order relative to the mean $M^2$.

Introduce the formal Laurent series at $s=\infty$
\begin{equation}\label{eq:dformal}
 d_\alpha(s)=\sum_{j\ge1}
       \frac{2\eta(2j)}{\alpha^{2j}}U_{2j-1}(s).
\end{equation}
The notation is asymptotic/formal: no convergence of the infinite sum in $j$
is asserted. Each coefficient of $s^{-1}$ uses only finitely many summands,
since $U_m(s)=O_m(s^{-m})$. The expansion begins
\begin{align}\label{eq:dseries}
 d_\alpha(s)={}&\frac c s+\frac c{s^2}+\frac c{s^3}
 +\frac{c+7c^2/5}{s^4}
 +\frac{c+77c^2/10}{s^5}\notag\\
 &+\frac{c+49c^2/2+372c^3/35}{s^6}+\cdots.
\end{align}
In the all-orders logarithmic sense,
\begin{equation}\label{eq:Fformal}
 F(t)\sim\alpha M\bigl(s/2-1+d_\alpha(s)\bigr).
\end{equation}
The lattice error $O(s)$ and the non-boundary continuum pieces are smaller
than $M/s^R$ for every fixed $R$. They are \emph{not} thereby zero and must
not be discarded when a multiplicative coefficient equivalent is wanted.

\section{Eliminating the saddle and generating all coefficients}\label{sec:eliminate}
Since $\dd t/\dd s=-\alpha(s-1)/M$, differentiating the continuum expansion
and using \eqref{eq:meanlattice} gives the formal mean relation
\begin{equation}\label{eq:Q}
 n\sim\frac{M^2}{2}Q_\alpha(s),\qquad
 Q_\alpha(s)=1+\frac{2(d_\alpha(s)+d_\alpha'(s))}{s-1}.
\end{equation}
Here $d_\alpha'$ denotes formal differentiation. Thus
\begin{equation}\label{eq:rfroms}
 r=s+\tfrac12\log Q_\alpha(s).
\end{equation}
Because $Q_\alpha(s)=1+O(s^{-2})$, this has a unique formal inverse
$s=r+O(r^{-2})$. It is computed recursively, for example by iterating
$s\leftarrow r-\frac12\log Q_\alpha(s)$ in the ring of formal series in
$r^{-1}$. Every iteration determines additional coefficients.

Substitution in $F(t)+nt$ gives
\begin{equation}\label{eq:formalH}
 \asym(r)=
 \frac{s(Q_\alpha(s)+1)/2-1+d_\alpha(s)}{\sqrt{Q_\alpha(s)}},
 \qquad r=s+\tfrac12\log Q_\alpha(s).
\end{equation}
By construction,
\begin{equation}\label{eq:Hseries}
 \asym(r)=r-1+\sum_{j\ge1}P_j(c)r^{-j}.
\end{equation}
Equations \eqref{eq:Urec}, \eqref{eq:dformal}, \eqref{eq:Q}, and
\eqref{eq:formalH} are an all-orders algorithm involving only rational
operations and expansions of logarithms and square roots. Since even zeta
values are rational multiples of powers of $\pi^2$, the coefficients are
polynomials over $\mathbb Q$ in $c$.

\begin{proof}[Proof of Theorem~\ref{thm:log}]
Fix the desired order $R$. Apply Lemma~\ref{lem:smoothing} and its mean
version with more than enough terms that both remainders, after all
operations, are $O_{\alpha,R}(r^{-R-1})$ in the normalized entropy.
The exact mean relation has an additional error $O(M)$, negligible compared
with $M^2r^{-A}$ for every fixed $A$. Inverting
\eqref{eq:rfroms} is stable because its derivative is $1+O(s^{-3})$.
Finite Taylor substitution is consequently legitimate with remainders, not
only as a formal manipulation.

We obtain
\[
 F(t_n)+nt_n=\alpha K\left(r-1+\sum_{j=1}^R P_j(c)r^{-j}
                         +O_{\alpha,R}(r^{-R-1})\right).
\]
Equation \eqref{eq:logsaddle} contributes only $O(s)$, which is
$o(Kr^{-R-1})$. This proves \eqref{eq:logall}. Direct expansion gives
\eqref{eq:P}; the supplied symbolic script reproduces those coefficients.
\end{proof}

\paragraph{One further coefficient for an independent check.}
The same algorithm gives
\begin{equation}\label{eq:P8}
 P_8(c)=c+\frac{259}2c^2+\frac{7271}{10}c^3
                  +\frac{157813}{840}c^4.
\end{equation}
The continuum quadrature script checks the approach of the remainder after
$P_7$, multiplied by $r^8$, to this polynomial. It does not replace the
preceding proof or evaluate integer coefficients at those enormous sizes.

\section{\texorpdfstring{Lambert-$W$}{Lambert-W} inversion with an error theorem}\label{sec:inverse}
The coefficients are strictly increasing for $n\ge1$. Indeed, increasing the
largest part of a distinct partition by one is an injection into distinct
partitions of $n+1$, and strictly increases each positive norm weight.
Thus the threshold index $N_\alpha(Y)$ in Theorem~\ref{thm:inverse} is
well-defined for large $Y$.

The dominant inversion problem is
\[
 Y=\alpha K(\log K-1).
\]
Writing $w=\log K-1$ gives $we^w=Y/(\alpha e)$, proving
\eqref{eq:Wlead}. In particular $\log K_0=R_0$.
Set $K=K_0E$. The all-orders formal inversion is the single equation
\begin{equation}\label{eq:Einverse}
 E\,\asym(R_0+\log E)=R_0-1,
 \qquad E=1+O(R_0^{-2}).
\end{equation}
The derivative with respect to $E$ is $R_0+O(1)$ at $E=1$, so successive
coefficients are uniquely determined. In terms of $c$, they are
\begin{align}\label{eq:inverseE}
 E={}&1-\frac c{R_0^2}-\frac c{R_0^3}
       +\frac{3c^2/2-c}{R_0^4}
       -\frac{c+2c^2/5}{R_0^5}\notag\\
 &-\frac{c+46c^2/5+5c^3/2}{R_0^6}+\cdots,
\end{align}
and hence
\begin{align}\label{eq:inverseE2}
 E^2={}&1-\frac{2c}{R_0^2}-\frac{2c}{R_0^3}
       +\frac{4c^2-2c}{R_0^4}
       +\frac{6c^2/5-2c}{R_0^5}\notag\\
 &-\frac{2c+77c^2/5+8c^3}{R_0^6}+\cdots.
\end{align}
The asymptotic index is $K_0^2E^2/2$.

To justify inversion of the sequence rather than just of a formal expression,
fix an order, use Theorem~\ref{thm:log} two orders farther than needed, and
bracket the target by the corresponding smooth functions with the two signs
of its error bound. Their derivatives with respect to $K$ are
$\alpha r(1+o(1))>0$. An entropy error $O(Kr^{-R-1})$ therefore changes $K$
by $O(Kr^{-R-2})$, and changes $n=K^2/2$ by
$O(K^2r^{-R-2})$. These bounds transfer the formal inversion to the actual
threshold by monotonicity. Rounding the brackets to integers costs $O(1)$,
smaller than all of these fixed-order error scales. This proves
Theorem~\ref{thm:inverse} and its all-orders extension.

\begin{remark}[Two different Lambert branches]
The saddle phase uses $W_{-1}$ because $s\to\infty$ solves
$se^{-s}=t/\alpha\to0^+$. The threshold inverse uses $W_0$ at a positive
argument. Interchanging the branches would destroy the relevant scale.
If inversion in elementary logarithms is desired, the usual formal solution
of $w+\log w=\log(Y/(\alpha e))$ can be substituted into
\eqref{eq:inverseE2}; retaining $W_0$ is shorter and preserves more precision.
\end{remark}

\section{Conditioned geometry: a logistic boundary and a Gumbel maximum}
Conditioning the Bernoulli law \eqref{eq:pk} on $S=n$ gives exactly
\begin{equation}\label{eq:conditioned}
 \Prob_n(\lambda)=\frac{\normp(\lambda)^\alpha}{a_\alpha(n)}.
\end{equation}
No approximation is involved in this identification.

\subsection{The boundary profile}
Fix distinct integer sites $k_{1,n},\ldots,k_{d,n}$ with
$(k_{j,n}-M)/H\to x_j\in\R$. Their unconditioned probabilities tend to
$p(x_j)=(1+e^{x_j})^{-1}$ by Lemma~\ref{lem:scales}.
For any prescribed pattern $\epsilon_j\in\{0,1\}$, independence gives
\begin{align}\label{eq:pattern}
 &\Prob_n(X_{k_{j,n}}=\epsilon_j\text{ for all }j)\notag\\
 &\quad=\prod_{j=1}^d p_{k_{j,n}}^{\epsilon_j}
                       (1-p_{k_{j,n}})^{1-\epsilon_j}
 \frac{\Prob_t\{S^{\rm rem}=n-\sum_jk_{j,n}\epsilon_j\}}
      {\Prob_t(S=n)}.
\end{align}
The shifted target differs from the remaining mean by $O(M)=o(\sqrt B)$.
The uniform local limit theorem in Section~\ref{sec:uniform} makes the ratio
of point probabilities tend to one. This proves both the logistic limit and
finite-site asymptotic independence. Similarly, a site at $\gamma M$ has
limiting occupation $1$ for $0<\gamma<1$, and $0$ for $\gamma>1$.

\subsection{The correct centering of the maximum}
The large solution of \eqref{eq:bdef} exists uniquely for sufficiently small
$t$. With $z=tb$, the equation is $z^\alpha e^{-z}=t^{\alpha+1}$, so
\begin{equation}\label{eq:blambert}
 b(t)=-\frac\alpha t\,
 W_{-1}\left(-\frac{t^{1+1/\alpha}}\alpha\right).
\end{equation}
Solving its logarithm,
$z-\alpha\log z=-(\alpha+1)\log t$, with $t=\alpha s e^{-s}$ yields
\begin{align}\label{eq:bexpand}
 \frac bM={}&1+\frac1\alpha
 +\frac{-\log s+\alpha\log(\alpha+1)-(\alpha+1)\log\alpha}
                {\alpha s}
 +O_\alpha\left(\frac{\log s}{s^2}\right).
\end{align}
Keeping only $b\sim(1+1/\alpha)M$ is not adequate for the Gumbel
centering: the omitted shift has size $M\log s/s$, larger than $1/t$.
Keeping all the displayed terms in \eqref{eq:bexpand} does give an
$o(1/t)$ centering error. The exact definition \eqref{eq:bdef} or
\eqref{eq:blambert} avoids this truncation bookkeeping altogether.

\begin{proof}[Proof of the maximum assertion in Theorem~\ref{thm:geometry}]
Fix $x\in\R$ and put $m=\lfloor b+x/t\rfloor$. Then $tm\to\infty$ and
$m/M\to1+1/\alpha>1$. In the tail $k>m$, all odds tend uniformly to zero.
A sum--integral comparison followed by integration by parts gives
\begin{align}\label{eq:tailint}
 \sum_{k>m}p_k
 &\sim\sum_{k>m}k^\alpha e^{-tk}
 \sim\frac{m^\alpha e^{-tm}}t
 \longrightarrow e^{-x}.
\end{align}
For the second equivalence, substitute $v=tx$ in the integral and use
$\int_z^\infty v^\alpha e^{-v}\dd v
=z^\alpha e^{-z}(1+O_\alpha(z^{-1}))$.
The endpoint discretization has relative size $O(t)$.
The last limit follows from $b^\alpha e^{-tb}/t=1$ and
$m/b\to1$. Also, $\sup_{k>m}p_k=O(t)$, so $\sum_{k>m}p_k^2=o(1)$.
Therefore the unconditioned probability of no part above $m$ is
\begin{equation}\label{eq:tailzero}
 \prod_{k>m}(1-p_k)\longrightarrow e^{-e^{-x}}.
\end{equation}

It remains to justify conditioning, rather than silently replacing one
ensemble by another. The removed mean and variance satisfy
\[
 \sum_{k>m}kp_k=O(M),\qquad
 \sum_{k>m}k^2p_k(1-p_k)=O(M^2).
\]
These bounds follow from the same tail-integral calculation with one or two
extra powers of $k$. The remaining sum $S_{\le m}$ thus has mean
$n-O(M)$ and variance $B(1+o(1))$. Its central block around $M$ is intact.
The uniform local limit estimate gives
\[
 \frac{\Prob_t(S_{\le m}=n)}{\Prob_t(S=n)}\longrightarrow1.
\]
Independence of the retained and removed sites now gives
\[
 \Prob_n(L_n\le m)=
       \prod_{k>m}(1-p_k)
       \frac{\Prob_t(S_{\le m}=n)}{\Prob_t(S=n)}
       \longrightarrow e^{-e^{-x}}.
\]
This proves the Gumbel limit. Its tightness, \eqref{eq:bexpand},
$M\sim K$, and $(tM)^{-1}=1/(\alpha s)\to0$ imply
$L_n/K\to1+1/\alpha$ in probability.
\end{proof}

\paragraph{The two locations are genuinely different.}
The occupation boundary lies near $M\sim\sqrt{2n}$, while the maximum lies
near $(1+1/\alpha)M$. For $\alpha=1$, a typical largest part is asymptotic
to $2\sqrt{2n}$, not to $\sqrt{2n}$. Individual occupancies far above the
boundary are small, but there are enough possible sites to move the extreme
by a macroscopic multiple of $M$. This is why a boundary profile alone does
not determine the largest part.

\section{A proved exponential scale from Fourier resonances}\label{sec:resonance}
The bound \eqref{eq:noncentral} suggests an exponentially small scale not
visible in either the logarithmic block or the central Edgeworth series.
We can identify its constant at each fixed integer resonance.

\begin{theorem}[Fixed-resonance modulus]\label{thm:resonance}
For every fixed nonzero integer $\ell$, with $\phi_t$ as in
Section~\ref{sec:llt},
\begin{equation}\label{eq:resonance}
 -\log\left|\phi_t\left(\frac{2\pi\ell}{M}\right)\right|
 \sim\frac{2\pi^4\ell^2}{3}\frac{H^3}{M^2}
 \sim\frac{2\pi^4\ell^2}{3\alpha^3}\frac M{s^3}.
\end{equation}
\end{theorem}
\begin{proof}
The exact modulus formula is
\[
 \log|\phi_t(\theta)|
  =\frac12\sum_{k\ge1}\log\left(1-4p_k(1-p_k)
                                    \sin^2(k\theta/2)\right).
\]
At $\theta=2\pi\ell/M$, write $k=M+xH$. For bounded $x$,
\[
 \sin^2(\pi\ell k/M)
 =\pi^2\ell^2 x^2H^2/M^2+O_\ell(x^4H^4/M^4).
\]
The logarithm is therefore $-4p_k(1-p_k)$ times this sine square to leading
order. Lemma~\ref{lem:scales}, applied with moments $2$ and $4$, gives
\[
 -\log|\phi_t(2\pi\ell/M)|
 \sim\frac{2\pi^2\ell^2H^3}{M^2}
     \int_\R x^2\frac{e^x}{(1+e^x)^2}\dd x.
\]
One may first restrict to $|x|\le s^\delta$ with small fixed $\delta>0$;
Taylor remainders are then uniformly controlled. Outside this interval,
the exponential boundary tails and the lower-endpoint bounds in
Lemma~\ref{lem:scales} are negligible at the displayed scale. The quadratic
logarithm error is $O(H^5/M^4)=o(H^3/M^2)$.
Finally,
\[
 \int_\R x^2\frac{e^x}{(1+e^x)^2}\dd x
 =4\sum_{j\ge1}\frac{(-1)^{j-1}}{j^2}=\frac{\pi^2}{3},
\]
by expansion on the positive half-line. This proves the theorem.
\end{proof}

For fixed $\alpha$, the scale $\exp[-C_\alpha M/s^3]$ is smaller than
every power of $s/M$, because $M/s^3$ dominates $\log M$. However, its
constant deteriorates as $\alpha^{-3}$. This makes the fixed-$\alpha$
asymptotic theorem compatible with large visible resonances at moderate $n$.

The theorem concerns the \emph{modulus at a specified point}. It does not
yet determine the shifted secondary saddles, their phases, or the contribution
of their full neighborhoods to the coefficient integral. Calling it a
complete multi-saddle transseries would be incorrect. It does identify a
rigorous exponential scale and a concrete next target for exponential
improvement.

\section{Reproducible computation and finite-size diagnostics}\label{sec:checks}
\subsection{Exact calculations}
The supplied \texttt{verify.py} computes the five integer sequences through
$n=5000$ by descending polynomial multiplication. It compares their initial
sixteen terms against the consulted OEIS entries and independently checks
$n\le150$ by a logarithmic-derivative recurrence. Explicitly, define
\begin{equation}\label{eq:exactB}
 C_\alpha(m)=\sum_{d\mid m}(-1)^{m/d+1}d^{1+\alpha m/d}.
\end{equation}
Then $a_\alpha(0)=1$ and
\begin{equation}\label{eq:exactrec}
 n a_\alpha(n)=\sum_{m=1}^n C_\alpha(m)a_\alpha(n-m).
\end{equation}
This follows by expanding $qA_\alpha'(q)/A_\alpha(q)$ formally. The two
algorithms have different arithmetic paths, so their agreement is a useful
implementation check. It is not a proof of the asymptotic theorems.

All five exact prefix checks, all recurrence checks through $150$, and strict
monotonicity checks through $5000$ passed. The complete computed lists are
included as plain text. They were independently generated here, not copied
from OEIS b-files.

\subsection{The central approximation is excellent for one column, not all}
Let $G=e^{F(t_n)+nt_n}/\sqrt{2\pi B}$, with the infinite sums evaluated by
high-precision truncation. The following are numerical diagnostics, not
interval-certified bounds. Values shown are rounded from the included CSV.
\begin{center}
\begin{tabular}{rrrrr}
\toprule
$\alpha$ & $n$ & $a_\alpha(n)/G$ & $a_\alpha(n)/(G(1+\mathcal E_1))$
 & $a_\alpha(n)/(G(1+\mathcal E_1+\mathcal E_2))$\\
\midrule
1 & 100 & 0.989159702 & 0.999826713 & 0.999908591\\
1 & 1000 & 0.997526322 & 0.999995495 & 0.999999984\\
1 & 5000 & 0.999084078 & 0.999999367 & 0.999999999\\
2 & 5000 & 0.999545155 & 1.000036354 & 1.000036527\\
3 & 5000 & 0.961625508 & 0.961944240 & 0.961944315\\
4 & 5000 & 1.230518103 & 1.230825332 & 1.230825386\\
5 & 5000 & 1.581097698 & 1.581414178 & 1.581414223\\
\bottomrule
\end{tabular}
\end{center}
At $\alpha=1,n=5000$, two central corrections reduce the relative error to
about $1.1\times10^{-9}$. At $\alpha=5,n=5000$, the central approximation
is plainly not accurate, and adding central corrections does not repair it.
This is evidence about finite-size behavior, not a contradiction of a
fixed-parameter limit. In particular, these tests must not be advertised as
numerical confirmation of rapid convergence for all five columns.

The resonance diagnostic computes the exact product modulus at
$\theta=2\pi/M$, using the same numerical saddle:
\begin{center}
\begin{tabular}{rrr}
\toprule
$\alpha$ & $n$ & $|\phi_t(2\pi/M)|$\\
\midrule
1 & 5000 & $2.33482\times10^{-15}$\\
2 & 5000 & $2.72621\times10^{-5}$\\
3 & 5000 & $0.0173352$\\
4 & 5000 & $0.150645$\\
5 & 5000 & $0.361805$\\
\bottomrule
\end{tabular}
\end{center}
The increase in these noncentral moduli is consistent with the observed
failure of a central-only approximation. Their values do not by themselves
predict the sign of the coefficient error; the omitted phases matter.

\subsection{Logarithmic expansions are not monotone numerical improvements}
For $\alpha=1,n=5000$, the exact logarithm is approximately
$401.567349464797$. Subtracting the leading expression $K(r-1)$ leaves
$41.05033086599$. Subtracting the expansion through $P_3$ instead leaves
$-2.72431695876$, while truncating through $P_7$ leaves $-7.67561845406$.
The additional terms are not a promise of monotone improvement at a fixed
small $r$. The missing prefactor and endpoint contributions are also on
scales suppressed by the logarithmic theorem but relevant numerically.

An independent continuum check evaluates the integral and mean in
Lemma~\ref{lem:smoothing}, without using the rational-function expansion,
and compares the normalized remainder after $P_7$ with $P_8$:
\begin{center}
\begin{tabular}{rrr}
\toprule
$\alpha$ & $s$ & Scaled remainder $r^8(\text{entropy}/(\alpha K)-\mathscr H_{\alpha,7}(r))$\\
\midrule
1 & 24 & 7881.08213\\
1 & 48 & 5997.09948\\
1 & 96 & 5424.47363\\
\multicolumn{2}{r}{Predicted limit $P_8(\pi^2/6)$:} & 4963.76905\\
3 & 24 & 15.7012553\\
3 & 48 & 9.85270737\\
3 & 96 & 9.49113831\\
\multicolumn{2}{r}{Predicted limit $P_8(\pi^2/54)$:} & 9.15764483\\
\bottomrule
\end{tabular}
\end{center}
The approach is slow, but has the expected direction in these samples.
These are continuum diagnostics at nonintegral implied $n$, not exact
integer-sequence computations at exponentially large indices.

\subsection{What the numerical implementation guarantees}
The integer arithmetic and recurrence comparisons are exact. For the saddle
calculations, the program uses arbitrary-precision arithmetic and chooses a
cutoff from a geometric majorant for
$\sum k^{\alpha+6}e^{-tk}$. This bounds the tails relevant to all cumulants
through order six, apart from harmless fixed constants. The test used
50 requested decimal digits with guard digits; the saved mean residuals are
far smaller than the displayed precision. Nevertheless, the majorant and
Newton iteration are evaluated in floating point, not outward-rounded
interval arithmetic. Accordingly, the CSV values are diagnostics rather than
certificates of rigorous numerical error bounds.

\section{Further research questions and proposed projects}\label{sec:future}
The proved results leave several substantial and concrete directions.

\paragraph{1. An effective multi-saddle expansion.}
Locate the secondary saddles near $2\pi\ell/M$, compute their phases and
Gaussian factors, and sum the relevant arcs with a uniform remainder. Theorem
\ref{thm:resonance} supplies the first exponential action but not the full
contribution. The moderate-$n$ failures for $\alpha=3,4,5$ make this a practical
as well as theoretical priority.

\paragraph{2. A complete two-scale coefficient transseries.}
Refine the continuum-to-lattice passage sufficiently to retain endpoint terms,
Euler--Maclaurin corrections, the central determinant, and the resonance
sectors in one consistent normalization. Determine which of these expansions
are Gevrey and whether optimal truncation produces a useful exponential
improvement. The present all-logarithmic series alone does not settle this.

\paragraph{3. The crossover as the tilt vanishes.}
The fixed-$\alpha$ expansion is singular as $\alpha\downarrow0$, while the
unweighted distinct-partition count has a different growth law. A natural
crossover parameter is $\alpha\log K$. Derive a uniform family of saddle
profiles when this parameter stays bounded, and match it to both the
unweighted regime and the fixed-positive-tilt regime. The shrinking-tilt
analysis in Gikunda's thesis provides an important comparison point
\cite{gikunda}, but matching the full regimes still requires work.

\paragraph{4. Growing moments and concentration near maximal norm.}
Allow $\alpha=\alpha(n)\to\infty$. The fluctuation width $H$ and resonance
action show that several thresholds can occur before the model collapses onto
maximal-norm partitions. Determine the sharp parameter ranges for a local
limit theorem, for single-saddle dominance, and for localization near an
extremal partition. The fixed-$\alpha$ proof does not justify substituting an
arbitrary growing parameter.

\paragraph{5. The joint law of the largest several parts.}
The independent tail calculation suggests a Poisson process with intensity
$e^{-x}\dd x$ for the rescaled upper parts. Prove its conditioned point-process
limit and obtain joint distributions of the largest $j$ parts, including
second-order corrections to their centering. The maximum theorem is the
one-dimensional starting point, not a substitute for tightness of the process.

\paragraph{6. Fluctuations of the number of parts and of the log norm.}
A multivariate conditioned local limit theorem should describe the number of
occupied sites and the logarithm of the norm under fixed positive tilt.
Conditioning removes a covariance direction and can cancel a leading
variance term; computing the surviving scale requires more than plugging the
unconditioned variance into a Gaussian formula.

\paragraph{7. More general weights.}
Replace $k^\alpha$ by $k^\alpha(\log k)^\beta$ away from the harmless small
indices, or by a smooth regularly varying weight. Identify hypotheses that
preserve a single large crossing, the logistic boundary, and a computable
inverse expansion. Arithmetic thinning of the allowed part sizes would also
require a replacement for the consecutive-block lemma and may introduce
persistent lattice obstructions.

\paragraph{8. Certified numerics and formal verification.}
Give explicit constants in the Fourier bounds, add interval-enclosed saddle
solving, and certify a chosen relative-error target at a stated $n$ and
$\alpha$. For Lean, a sensible modular route is: the formal Euler product and
recurrence; the Bernoulli coefficient identity; the consecutive-block
inequality; an abstract triangular-array local-limit theorem; and finally the
boundary smoothing lemma. The algebraic generation of $P_j$ and the inverse
coefficients can be separately certified by finite polynomial identities.
This would connect directly to the repository's transseries infrastructure
without pretending that formal algebra alone proves analytic asymptotics.

\section{Conclusion}
The five checked OEIS conjectures share a single analytic mechanism. Filling
parts up to a moving boundary produces the dominant factorial-scale entropy;
smoothing that boundary produces the $\pi^2$ correction and its all-orders
inverse-logarithmic successors. A conditioned Bernoulli model supplies the
local probability that converts this entropy into an actual coefficient,
and also gives logistic occupation and a Gumbel extreme. Lambert-$W$
inversion turns the same expansion into a quantitative threshold law.

The resulting package contains explicit proofs rather than numerical
extrapolation. It also records an important limitation: central asymptotics
can converge extremely slowly for higher moments because noncentral
resonances are only suppressed on the scale $\exp[-C_\alpha K/(\log K)^3]$.
That proved scale and the visible finite-size discrepancies identify a
well-defined next research problem: a phase-sensitive, exponentially
improved multi-saddle expansion.

\appendix
\section{Reproduction guide and file manifest}
The archive contains this source and its compiled PDF, together with:
\begin{center}
\begin{tabular}{p{.38\linewidth}p{.53\linewidth}}
\toprule
File & Purpose\\
\midrule
\texttt{verify.py} & Exact integer sequences, independent recurrence, and numerical saddle/Edgeworth diagnostics.\\
\texttt{derive\_series.py} & Symbolic generation of the boundary derivatives and logarithmic coefficients.\\
\texttt{derive\_inverse.py} & Symbolic reversion through six inverse-logarithmic orders.\\
\texttt{check\_continuum.py} & Independent continuum quadrature diagnostics.\\
\texttt{check\_resonances.py} & Direct numerical product moduli at the first resonance.\\
\texttt{data/} & Computed integer lists, CSV diagnostics, symbolic output, and execution logs.\\
\texttt{README.md}, \texttt{requirements.txt} & Build and reproduction instructions, scope, and dependencies.\\
\bottomrule
\end{tabular}
\end{center}
The mathematical proofs do not require running the programs. To reproduce:\par\medskip
\noindent\begin{minipage}{\linewidth}
\begin{lstlisting}
python -m pip install -r requirements.txt
python verify.py --max-n 5000 --dps 50
python derive_series.py
python derive_inverse.py
python check_continuum.py
python check_resonances.py
pdflatex -interaction=nonstopmode -halt-on-error article.tex
pdflatex -interaction=nonstopmode -halt-on-error article.tex
\end{lstlisting}
\end{minipage}\par\medskip
The source uses standard LaTeX packages and no external bibliography processor.
No Lean code or machine-checked proof is bundled or implied.

\section{Compact proposed OEIS formula updates}
For the entry corresponding to $\alpha=1,\ldots,5$, write
$K=\sqrt{2n}$ and $r=\log K$. The following common formula is proved in
Theorem~\ref{thm:log}:
\[
 \log a_\alpha(n)=\alpha K(r-1)
       +\frac{\pi^2K}{6\alpha r}
       +\frac{\pi^2K}{6\alpha r^2}
       +O_\alpha(K/r^3).
\]
It is appropriate to link a reviewed proof before changing a database's
conjectural status. This report has not submitted or edited any OEIS entry.
For the array \OEIS{A292189}, the result holds down every fixed positive
column; it makes no uniform assertion about the main diagonal or columns
whose index grows with the row.

\clearpage
{\small\raggedright
\begin{thebibliography}{99}
\bibitem{oeis1} OEIS Foundation Inc., \emph{A022629: Expansion of
$\prod_{m\ge1}(1+mq^m)$}. Conjectural formula attributed to V. Kote\v{s}ovec,
May 8, 2018. Internal entry consulted October 1, 2026; sequence revision
\#54, March 13, 2026. \url{https://oeis.org/A022629/internal}.
\bibitem{oeis2} OEIS Foundation Inc., \emph{A092484: Expansion of
$\prod_{m\ge1}(1+m^2q^m)$}. Conjectural formula attributed to V. Kote\v{s}ovec,
December 27, 2020. Consulted October 1, 2026.
\url{https://oeis.org/A092484}.
\bibitem{oeis3} OEIS Foundation Inc., \emph{A265840: Expansion of
$\prod_{k\ge1}(1+k^3q^k)$}. Conjectural formula attributed to V. Kote\v{s}ovec,
December 27, 2020. Consulted October 1, 2026.
\url{https://oeis.org/A265840}.
\bibitem{oeis4} OEIS Foundation Inc., \emph{A265841: Expansion of
$\prod_{k\ge1}(1+k^4q^k)$}. Conjectural formula attributed to V. Kote\v{s}ovec,
December 27, 2020. Consulted October 1, 2026.
\url{https://oeis.org/A265841}.
\bibitem{oeis5} OEIS Foundation Inc., \emph{A265842: Expansion of
$\prod_{k\ge1}(1+k^5q^k)$}. Conjectural formula attributed to V. Kote\v{s}ovec,
December 27, 2020. Consulted October 1, 2026.
\url{https://oeis.org/A265842}.
\bibitem{oeisarray} OEIS Foundation Inc., \emph{A292189: Square array of
coefficients of $\prod_{j\ge1}(1+j^kq^j)$}. Consulted October 1, 2026.
\url{https://oeis.org/A292189}.
\bibitem{schneider} R. Schneider and A. V. Sills,
\emph{The product of parts or ``norm'' of a partition}, Integers
\textbf{20A} (2020), Article A13. Author version:
\url{https://arxiv.org/html/1904.08004v2}.
\bibitem{bridges} W. Bridges and W. Craig,
\emph{On the distribution of the norm of partitions}, The Ramanujan Journal
\textbf{66} (2025), Article 40. DOI:
\href{https://doi.org/10.1007/s11139-024-00969-5}{10.1007/s11139-024-00969-5}.
Author version: \url{https://arxiv.org/html/2308.00123v1}.
\bibitem{rana} M. Rana, H. Kaur, and A. Kumar,
\emph{On the sum of partition norms and its connection to norms of partitions
with parts greater than one}, Notes on Number Theory and Discrete Mathematics
\textbf{31} (2025), no. 4, 908--915. DOI:
\href{https://doi.org/10.7546/nntdm.2025.31.4.908-915}{10.7546/nntdm.2025.31.4.908-915}.
\url{https://nntdm.net/papers/nntdm-31/NNTDM-31-4-908-915.pdf}.
\bibitem{gikunda} D. K. Gikunda,
\emph{The Distribution of the Product of Parts in Integer Partitions},
Ph.D. dissertation, Stellenbosch University, March 2026, Chapter 4,
especially Section 4.2 and Proposition 4.2.5. Author-uploaded text consulted
October 1, 2026. DOI:
\href{https://doi.org/10.13140/RG.2.2.22314.07368}{10.13140/RG.2.2.22314.07368}.
Repository record: \url{https://scholar.sun.ac.za/handle/10019.1/136012}.
\bibitem{proveit} V. Reshetnikov,
\emph{ProveIt}, repository snapshot
\texttt{1f1981f682b2878bde51a6ad40c22777f362fc05},
\texttt{Analysis/Transseries/README.md}, consulted October 1, 2026.
\url{https://github.com/VladimirReshetnikov/ProveIt}.
\end{thebibliography}
}
\end{document}
